\chapter{Introducción}
    El desarrollo de proyectos arquitectónicos, concretamente en el aspecto de la visualización, ha evolucionado de manera constante desde los primeros planos dibujados a mano hasta los modelos digitales tridimensionales de hoy en día. Durante décadas, arquitectos e ingenieros han confiado en el dibujo técnico de representaciones bidimensionales, planos, secciones y alzados para explicar sus diseños a constructores, inversores y clientes. Sin embargo, esta forma de comunicación tiene limitaciones, ya que requiere de conocimientos técnicos específicos para su correcta interpretación y hace difícil la comprensión espacial completa del proyecto, especialmente para personas sin formación en el sector.

    Con la aparición de la metodología BIM (Building Information Modeling), esta industria dio un salto considerable al juntar geometría tridimensional con información técnica detallada de cada elemento de la construcción. No obstante, incluso los modelos BIM más avanzados se visualizan habitualmente en pantallas convencionales, lo que mantiene esa dificultad perceptiva entre el diseño digital y la experiencia física del espacio construido. Esto hace que se generen problemas de manera recurrente: malinterpretaciones o dificultades para validar decisiones durante la fase de ejecución y barreras comunicativas con clientes.

    De manera paralela, en la última década, las tecnologías de Realidad Aumentada (AR), y Realidad Virtual (VR) se han desarrollado o perfeccionado lo suficiente para que sean aplicables en estos sectores de manera profesional. Dispositivos como Microsoft HoloLens, Meta Quest o smartphones de gama media ya poseen capacidades de seguimiento espacial, renderizado en tiempo real y experiencias inmersivas a costes cada vez más asequibles. Estas tecnologías permiten superponer información digital sobre el entorno físico real (AR) o sumergir completamente al usuario en un espacio virtual (VR), abriendo posibilidades para el sector de la construcción.

    Actualmente, diversas empresas y estudios de arquitectura han comenzado a experimentar con soluciones AR/VR de forma aislada, por ejemplo, visualizar el interior de un piso para su compra. Sin embargo, la mayoría son herramientas propietarias, es decir, carecen de integración fluida con flujos de trabajo BIM existentes o están diseñadas para casos de uso muy específicos. Por lo tanto, existe la posibilidad de desarrollar una aplicación integral que aproveche estas tecnologías inmersivas y satisfaga las necesidades reales de los profesionales del sector.

    El objetivo de este trabajo es diseñar y desarrollar una aplicación basada en Realidad Virtual y/o Realidad Aumentada orientada específicamente al ámbito de la construcción de edificios. El objetivo de este proyecto permitirá visualizar modelos BIM directamente sobre el emplazamiento real mediante dispositivos AR, facilitando la inspección de obra, la detección temprana de errores entre proyecto y ejecución, y la comunicación efectiva entre los equipos de las distintas fases. Asimismo, ofrecerá experiencias inmersivas en VR para recorrer virtualmente el proyecto antes de su construcción, evaluar su diseño, realizar formaciones en procedimientos de seguridad y presentar propuestas a clientes de manera comprensible e impactantes.

    Este proyecto se estructura en varias fases. En primer lugar, se llevará a cabo un análisis exhaustivo del estado del arte en tecnologías en AR/VR aplicadas a la construcción, identificando las soluciones existentes, sus limitaciones y las necesidades no cubiertas. Después, se definirá la arquitectura del sistema, seleccionando las plataformas de desarrollo más adecuadas, estableciendo la integración con formatos BIM estándar (como IFC o Revit) y diseñando la interfaz de usuario para que su interacción sea intuitiva con entornos inmersivos. La fase de desarrollo incluirá la implementación de funcionalidades clave como la medición métrica, la generación de informes y la visualización de datos técnicos. Finalmente, se realizarán pruebas en edificios reales para validar la usabilidad y efectividad de la solución.

    Las tecnologías inmersivas eliminan la abstracción que suponen los planos tradicionales y acercan el diseño digital a la percepción humana del espacio real. Este trabajo busca contribuir a estas tecnologías, proporcionando una herramienta práctica que mejore la eficiencia, reduzca errores y potencie la colaboración en los proyectos de construcción.

    El resto de este documento está estructurado de la siguiente manera:
    
    En primer lugar, en el Capítulo 2 se presenta el estudio de viabilidad, donde se analiza el contexto y mercado del proyecto mediante herramientas: una matriz DAFO (Cuadro~\ref{tab:dafo}), un modelo Lean Canvas (Cuadro~\ref{tab:lean-canvas}) y un análisis de riesgos (Cuadro~\ref{tab:riesgos}).
    
    Posteriormente, el Capítulo 3 define los objetivos principales y específicos, incluyendo la planificación temporal y una revisión exhaustiva del estado del arte de las tecnologías inmersivas y BIM. 
    
    En el Capítulo 4 se expone la metodología de trabajo adoptada para el proyecto, basada en Kanban, así como las herramientas utilizadas para su gestión.
    
    Por otro lado, el Capítulo 5 entra en detalle sobre el desarrollo técnico de la aplicación, describiendo la implementación de los diferentes sistemas como la selección de elementos BIM.
    
    Finalmente, los Capítulos 6 y 7 recogen, respectivamente, los resultados obtenidos tras las pruebas de la aplicación y las conclusiones finales a partir de la realización de este trabajo.


